\section{Introduction}
In a busy world where hands-free interaction with devices is becoming increasingly important, voice-controlled robots offer a promising solution for various applications, from home automation to assistive technologies. This report presents the development of a voice-controlled robot using a TurtleBot3 platform, integrating speech recognition and sound localisation to enable intuitive user interaction. The robot will then be tested in varying noise conditions to measure its performance in noisy environments typical for silent, casual or workshop settings.

The primary objective of this project is to design and implement a system that allows users to control a mobile robot's movements through voice commands. The burden of interpretation is placed on the robot, allowing the user to simply say the commands, as opposed to a remote control device. For this, both artificial intelligence and signal processing techniques are leveraged.

Furthermore, the robot should be capable of locating the sound source of the voice commands. This will enable the user to recall the robot to their location by simply calling out to it, further enhancing the interactivity and usability of the system. This requires sound localisation capabilities, which are built in this project. %This system could for example be implemented in workshop environments for retrieving or moving items. The goal is therefore to create a voice-controlled robot than can be used hands-free, letting the user do other task with tkeir hands while still beng able to control the robot.

A way to get data from the microphones, transcribe the audio, and analyze the transcription for any command is required, and if the command found is a ``come here'' command, the robot should then find the person speaking and move towards them. For ease of use the robot should be able to identify both the command and location from a single audio recording, so the user doesn't have to issue multiple commands.

The robot should follow a layered software architecture, allowing for easy maintenance and scalability. Furthermore, a communication protocol also has be included to communicate between the different components.

The work distribution among the team members is shown in the appendix in section \ref{app:work-distribution}, by main contributer and supporting. The github repository for the project can be found in section \ref{app:github-repo}.